\head{第13节：指定点截面、Bertini归纳与正维中心}
\source{法稿物理78--79页及80页顶部NB；英编稿51页中部至52页NB末；旧注D17-M061--080}

\audit{V21-BER-E01}{源页恢复与错误归属}{逐页目视法稿，并以英文编订稿和旧中文逐段对照}
本模块从“现在只剩下证明引理13.3”开始，止于13.4标题之前。法稿78、79页为连续正文；80页只取顶部NB。英文位置为51页中部至52页NB末。已查看法稿78--80页及英文51、52页，并对原字母和删划另作高分辨率裁图检查；法77页仅作前一模块边界确认，法80页13.4以下不计入本模块。

本次没有以新的数学论证代替法文辨读。能够直接从法稿恢复的内容包括：任意给定射影浸入的限定、小写$x\in P(k)$、一般点下标$\eta$、两处$\varphi_n f$、表示“问题”的缩写Pb、固定$x$而超平面一般的语序、$f(X)$、先由$H$再由$L'$截取的完整归纳步骤，以及两处正确的条件连接词。法79页“正式表述中的繁重细节”也是完整句子。英文中的$p^0$、$Q_n/\Phi_n$、空方括号及\emph{under/unless}等不再误记为原作者的缺文或逻辑错误。

但法稿确实把一般纤维写成“在$k$上光滑”；本版维持数学订正，改成在$\kappa(\eta)$上光滑。法78页的$\underline Y^{\mathrm{sing}}$与79页的$\underline X^{\mathrm{sing}}$也确有字母差异，不能仅因英文在前处写$h^{\mathrm{sing}}$就臆造一致的法稿。旧8.7(c)、8.8的缺续没有被这几页补全。

法80页“没有什么能排除$T$包含直线$C^0$”完整可读，故旧中文的连接语残缺标记撤销。$C^0$解释为射影对偶的消灭子参数空间；原文例取$\dim C=r-2$时，它恰是一条直线。冒号后两行有密集覆打删除，可辨的零碎字母不足以重建完整旧句；保留删稿的图像和位置，不将旁边未删文字倒填入删除层。

\audit{V21-BER-E02}{原证明的逻辑及依赖}{恢复正确的法文顺序，不将引用路线冒作已补齐的原文}
原稿证明$m=1$情形时，以闭坏参数集$T\subset P^\vee$和超平面$Q=x^0$比较维数。若$T$余维至少2，则$Q$的一般点不属于$T$。若$T$为不可约的余维1闭集，则$Q$的一般点属于$T$等价于$Q=T$。这两步的集合论与维数论证成立，法稿连接词没有错误。

随后的论证另用8.7(c)的对偶切空间结论，以及相应奇异关联态射一般平展的条件。若该对偶结论确实适用，$T=Q=x^0$便使一般$z$满足$f(z)=x$；因为$f^{-1}(\{x\})$是闭集并含稠密子集，从而$f(X)=\{x\}$。在有限型无分歧背景下，$f$拟有限，这迫使$X$零维。这里每项条件都有独立作用，不能删去无分歧性后把常值态射的定义域一概判为零维。

法稿的一般$m$归纳先由$H$截取，应用刚得到的$m=1$结论，再由$L'$截取，应用$m-1$的归纳假设。参数是在经过$x$的超平面及余维$m-1$子空间的乘积中独立一般选择；取交给出一般余维$m$子空间，允许相应的基域扩张。旧中文根据英文残缺句形成的反向顺序不再沿用。以下补证另外建立一个足以用于相应构造的Veronese版本，不宣称恢复了8.8原句，也不把特征零下不作Veronese的全部原断言包括在已独立重证范围内。

\audit{V21-BER-E03}{经过一个固定有理点的一般截面}{一阶截面与行列式余维的独立充分补证}
以下命题及证明为本轮译者补证。设$k$为任意域，$f\colon X\to\mathbf P^r_k$是有限型无分歧态射，$X$光滑且为$k$上有限型。取$x\in\mathbf P^r(k)$、整数$d\geq2$，并令
\[
\nu_d\colon\mathbf P^r\hookrightarrow\mathbf P^N,
\qquad N=\binom{r+d}{d}-1.
\]
设$1\leq m\leq N$，并假设$x$不属于$X$中任何维数小于$m$的不可约分量的像。则存在参数化包含$\nu_d(x)$的余维$m$线性子空间的格拉斯曼簇中的非空开集，使其中每个几何参数所给出的截面光滑。对于$X$的维数为$n\geq m$的几何分量，非空截面纯维数为$n-m$；维数小于$m$的分量与一般截面不交。这个命题不要求$X$固有，也不要求$k$无限；后一个条件只用于保证非空开参数集有$k$-有理点。

\begin{analysisproof}
置
\[
W=H^0(\mathbf P^r,\mathcal I_x(d)),\qquad
e=\dim_kW=N,
\]
并在仿射参数空间$A=W^m$中取有序$m$元截面$(s_1,\ldots,s_m)$。记$M=me$。维数计算可在代数闭包上进行。无分歧态射的几何纤维是有限的约化零维概形，因此$f^{-1}(x)$只有有限多个几何点。$X$的光滑性使其几何不可约分量互不相交，每一分量维数恒定；有限型性保证分量个数有限。

\lead{第一步：离开固定点的一阶满射。}
对几何点$y\ne x$，选择线性形式$\ell$，使$\ell(x)=0$、$\ell(y)\ne0$；另取$h$使$h(y)\ne0$。当$q$遍历所有线性形式时，
\[
\ell h^{d-2}q\in W
\]
在$y$处的值和一阶微分张成$\mathcal O(d)$的一阶截面空间。这是因为以$h^d$平凡化后，这些函数为$(\ell/h)(q/h)$；$\ell/h$在$y$处为单位，而$q/h$给出常数及全部线性坐标。对$f(z)=y$，无分歧性给出满射
\[
f^*\Omega_{\mathbf P^r/k}\twoheadrightarrow\Omega_{X/k}.
\]
故在$X$的每个这样的$z$处，$W$同样满射到截面在$z$处的值与全部一阶微分。

\lead{第二步：基点处的微分满射。}
在$x$处取$h(x)\ne0$，让$q$遍历在$x$消失的线性形式。截面$h^{d-1}q$属于$W$，并给出$x$处全部余切方向。再由上述余切满射，它们在每个$z\in f^{-1}(x)$处给出$X$的全部余切方向。所有这些截面在基点处的值均为零，这是本步与第一步的区别。

\lead{第三步：远离基点的坏关联集。}
固定$X$的一个维数为$n\geq m$的几何分量。在$z\notin f^{-1}(x)$处，要求
\[
s_1(z)=\cdots=s_m(z)=0
\]
给出$m$个独立线性条件。此后，$m$个微分构成一个$m\times n$矩阵。其秩小于$m$的条件具有余维$n-m+1$：秩至多$m-1$的矩阵簇由秩恰为$m-1$的稠密层控制，该层维数为$(m-1)(n+1)$，所以在维数$mn$的矩阵空间中余维为$n-m+1$。因此“同时消失而微分秩不足”的参数在固定$z$上总余维为
\[
m+(n-m+1)=n+1.
\]
令$z$变动后，相应坏关联集的维数至多为
\[
n+M-(n+1)=M-1.
\]
这里一阶满射可局部写成向量丛满射，故逐纤维余维计算确实给出关联集的维数界；不是仅作闭点个数计算。有限型态射的像的闭包维数不超过源维数，因而坏参数的闭包是$A$的真闭集。即使$X$不固有，也只需取像的闭包，不曾假定像本来闭。

\lead{第四步：基点和低维分量。}
若$z\in f^{-1}(x)$位于维数$n$的分量上，假设保证$n\geq m$。第二步表明微分矩阵遍历所有$m\times n$矩阵；秩不足因而是一个真闭条件。只有有限多个这样的几何点，其坏条件的并仍为真闭集。

若某一分量维数$n<m$，则按假设它不含基点。在其任一点上，$m$个值同时为零的条件已有余维$m$，因此相应零点关联集维数至多为$M+n-m<M$。一般参数便与整个分量不交，无须把一个负的“预期维数”误写成非空概形的维数。

将所有几何分量的坏集并起来，仍得到维数小于$M$的闭包。为使开集明确在$k$上定义，按光滑态射的相对维数，将$X$分解成有限多个在$k$上定义的开闭子概形$X_n$；在每个$X_n$上分别作上述值与微分秩的关联集，再取其像的闭包之并。几何基变换后的维数估计表明这个并是真闭集，故其补集给出$A$中的非空开集$U$。对于$U$中的任意几何参数，零点处的$m$个相对微分线性无关；在光滑的相对环境中，Jacobian判据遂给出相对维数$n-m$的光滑零点概形。这一步直接使用相对Jacobian条件，不把“某个纤维的集合上没有坏点”无条件等同于族的光滑性。

\lead{第五步：由有序方程下降到格拉斯曼参数。}
取$W^m$中线性无关的$m$元组组成的非空开集$F$。映射
\[
F\longrightarrow\operatorname{Gr}(m,W)
\]
是秩$m$泛子丛的标架丛，局部平凡，结构群为$\mathrm{GL}_m$，因此是光滑、满射且开映射。可把前述坏集及其闭包选成$\mathrm{GL}_m$不变的：可逆行变换不改变共同零点理想，也不改变零点处微分的秩。因此$U\cap F$为不变的非空开集，其像是格拉斯曼簇中的非空开集。

一个$m$维子空间$V\subset W$就是$m$个线性独立的、经过$\nu_d(x)$的超平面方程，故给出$\mathbf P^N$中包含该点的余维$m$子空间。改换方程基底不改变截面，光滑性又可沿域扩张忠实下降，所以上述性质确实属于格拉斯曼参数，而非某一任选标架。命题得证。
\end{analysisproof}

这个证明只用一阶截面分离、无分歧态射的余切满射、矩阵秩层维数及有限型像的维数界。其维数方法可与Stacks Project的Bertini证明比较：\url{https://stacks.math.columbia.edu/tag/0FD4}，特别是引理33.47.3；这里针对固定点和任意余维的具体计算已经另行写出。它没有证明截面的几何连通性；后者在13.2的初等纤维化构造中是另一项输入，不能由本命题的光滑性结论偷换而来。

\audit{V21-BER-E04}{原附记的二次曲面例}{固定正维中心确实可能使全部截面非光滑}
下面给出原稿“任意特征下取非奇异二次超曲面”这一提示的具体实现；这些坐标不冒充原稿已有的坐标。取任意代数闭域上的光滑二次曲面
\[
X=V(x_0x_3-x_1x_2)\subset\mathbf P^3,
\qquad C=V(x_0,x_1)\subset X.
\]
四个偏导数为$x_3,-x_2,-x_1,x_0$，不可能在射影点同时为零，故曲面在任意特征下都光滑。经过$C$的平面组成一条对偶直线，方程为
\[
a x_0+b x_1=0,
\qquad [a:b]\in\mathbf P^1.
\]
在$b\ne0$的图上消去$x_1$后，截面方程为
\[
x_0\bigl(x_3+(a/b)x_2\bigr)=0.
\]
它是两条相交直线，交点处不光滑。$b=0$时同样得到$x_1x_2=0$。因此整个参数直线$C^0$都在坏参数集中。这验证原NB的警示，并消除旧英文残句造成的方向疑问。

\audit{V21-BER-E05}{Veronese之后仍可有全坏的固定正维中心}{给原疑问的无限制肯定答案提供反例}
原文只是提出“对于$\varphi_d f$、$d\geq2$，这种现象或许不再发生”的疑问，不能把它报告成作者已经证明的一条错误定理。但在没有额外限制$C$的情况下，肯定答案确实不成立。

\begin{analysisproof}
在特征零代数闭域上，取$f=\id_{\mathbf P^1}$和三次Veronese嵌入
\[
\nu_3\colon\mathbf P^1\longrightarrow\mathbf P^3,
\qquad [u:v]\longmapsto[u^3:u^2v:uv^2:v^3].
\]
令$C=V(x_2,x_3)$。在$u\ne0$的图上，曲线参数为$(1,t,t^2,t^3)$，故$C$是$t=0$处的射影切线。其维数为1，等于目标空间$\mathbf P^3$的维数减2，因而经过它的平面确实组成一维族，并不是只有一个平面的退化情形。

包含$C$的平面方程为$a x_2+b x_3=0$，拉回后是
\[
a u v^2+b v^3=v^2(au+bv)=0.
\]
当$a\ne0$时，$v=0$为重数2的点，而$au+bv=0$给出另一个不同的单点。在$t=v/u$坐标下，第一点的局部环为$k[t]_{(t)}/(t^2)$，所以不光滑。当$a=0$时，截面是$v^3=0$，为一个长度3的点，也不光滑。每个平面成员都不光滑，因此固定中心参数空间整个落入坏参数集。
\end{analysisproof}

这已满足Veronese次数至少2的要求，而且$f$是闭浸入、特别无分歧。反例中的$C$与像相交并含该处切线；所以它只否定原段没有限制$C$时的普遍肯定答案，不能被转述为否定另加$C\cap(\nu_d\circ f)(X)=\varnothing$、横截性或其他条件的版本。正文以明确的译校订正代替无依据的乐观推测，原问句完整保留在脚注V21-BER014。

\audit{V21-BER-E06}{旧注的保存与本次状态}{旧译校意见作为历史保留，不让过时归责继续支配正文}
旧D17-M061--080二十条注的原文已经逐条提取到本模块的历史记录，并与继承文件中的字节核对。它们并未被悄悄删除；本版只是以实际法稿对读结果更新了它们在当前正文中的位置和效力。

\begin{longtable}{p{.22\textwidth}p{.66\textwidth}}
\toprule 旧注 & 本次处理\\\midrule\endhead
D17-M061--063 & 一般点、指定点和编写建议的范围保留；恢复法稿手添的“任意给定射影浸入”；大写$X$误排归于英文。\\
D17-M064--065 & $x\in P(k)$、下标$\eta$及$\varphi_d$按法稿恢复；不可约分量问号改标英编者；基域$\kappa(\eta)$订正保留且确认法稿也有原误。\\
D17-M066--069 & 恢复Pb、固定点语序及$f(X)$；记录法稿$\underline Y^{\mathrm{sing}}/\underline X^{\mathrm{sing}}$差异，旧英文$h$不沿用；未补齐8.8。\\
D17-M070--072 & 法稿完整的先$H$后$L'$归纳及末句全部恢复；撤销针对原稿语法和缺词的旧归责。\\
D17-M073--077 & 闭性、维数及对偶依赖继续保留；法稿\emph{que si}、\emph{lorsque}和推论链原本正确，英文误译据实改正。\\
D17-M078--080 & 正维中心警示恢复完整意思；两行密集删划仍不猜补；原Veronese疑问保留，并加入三次Veronese切线反例。\\
\bottomrule
\end{longtable}

\lead{仍未独立完成的事项。}
原8.7(c)、8.8的全部一般版本和缺续没有恢复；特征零下不改变给定浸入的一般固定点结论未在本模块另作完整证明。非有理闭点集合的共同选择与下降由本版13.4(c)核查另行处理，本节单点补证不假称已经同时处理所有闭点。几何连通性是另一项几何输入。作者删去的密集覆打旧文仍有不可辨部分。技术编译成功不改变这些边界。
